\chapter{Benchmarks: the public sets and the implant bench (contrast z, foreignness, PSF)}\label{ch:benchmarks}

\begin{onething}
No public labelled anomaly set tests a model that reads satellite reflectance by wavelength, so the repository puts
synthetic objects into its own held-out satellite scenes - mixed in linearly, blurred by a Gaussian point-spread
function, and solved so that their contrast $z$ and foreignness match real labelled objects \emph{as they are
realised} - and trusts the result only once it reproduces an ordering the real labels already showed.
\end{onething}

\section{Why a bigger benchmark was needed}\label{sec:benchmarks:why}

\Cref{ch:measures} gave you measures that look at the top of the ranking. On the four labelled PRISMA scenes
(909 objects\hfil\penalty50\hfilneg\tagM{research/experiments/implant-bench/README.md}) they still could not separate the models: every
measure tied the shipped model with the fine-tuned ones within scene-to-scene noise
(\file{research/experiments/indradhanu-finetune/FT02b-mixed/README.md}). Four scenes and 909 objects are too few to
see differences of the size being chased; \cref{ch:bootstrap} puts a number on ``too few''.

The obvious fix is to borrow a bigger public benchmark. On 2026-10-16 the two standard ones were downloaded and
vetted, and the record is \file{research/experiments/benchmarks/README.md}. This chapter first shows why neither
fits, then teaches the benchmark the repository built instead, and the mathematics of each of its parts:
linear mixing, the ring and its contrast $z$, foreignness, the Gaussian point-spread function and the strength solve.

\begin{analogy}
To test proofreaders, a publisher takes real pages and plants typos of known difficulty, some glaring, some a single
swapped letter. Real pages keep the surrounding text honest; planted errors make every answer known. The test is only
fair if the planted typos look like the ones writers actually make - so you study a pile of real typos first, and
you check that the planted test ranks last year's proofreaders in the order their real work did.
\end{analogy}

\section{What a model needs from each scene}\label{sec:benchmarks:needs}

The Indradhanu models read \emph{any band list by wavelength} and expect \emph{reflectance} on a 0--1 scale
(\cref{ch:light}). So every benchmark scene must supply four things:
\begin{enumerate}
\item the cube;
\item the centre wavelength of every band;
\item a ground-truth mask;
\item the value scale: what a stored number means physically.
\end{enumerate}
Miss any one and the model is either guessing or reading something it was never trained on. The evaluator enforces
the fourth: \code{Evaluation.run} raises \code{ValueError} for any scene whose \code{units} are not
\code{"reflectance"} (\file{app/evaluation/runner.py}).

\section{HAD100}\label{sec:benchmarks:had}

HAD100 was published with the AETNet paper \citep{li2023you} and is cut from the public AVIRIS-NG archive (the
airborne sensor of \cref{ch:sensors}). The numbers below are from the benchmarks README, section 2\hfil\penalty50\hfilneg\tagM{research/experiments/benchmarks/README.md}.

\begin{center}
\small
\begin{tabular}{p{2.8cm}p{9.6cm}}
\toprule
item & value \\
\midrule
what ships & 94 test tiles, their ground truth, and 522 anomaly-free training tiles \\
protocol & \texttt{main.py} crops 100 test scenes of 64 $\times$ 64 from the 94 tiles (6 tiles give two crops) \\
labels & 3,567 anomaly pixels in 460 objects (8-connected); per scene median 25.5 pixels (0.6\%) \\
values & L1B at-sensor \textbf{radiance}, $\mu$W\,nm$^{-1}$\,cm$^{-2}$\,sr$^{-1}$, not reflectance \\
wavelengths & 425 per tile, 376.7--2500.4 nm; they differ slightly between flight lines \\
\bottomrule
\end{tabular}
\end{center}

\index{radiance}\emph{Radiance} is the light arriving at the sensor, atmosphere included. \index{reflectance}\emph{Reflectance}
is the fraction of sunlight the ground sends back, with the sun and atmosphere taken out. The crudest bridge is
\index{top-of-atmosphere reflectance}\emph{top-of-atmosphere reflectance}:
\begin{equation}\label{eq:benchmarks:toa}
\rho_{\text{TOA}} = \frac{\pi L}{E_0 \cos\theta_s},
\end{equation}
with $L$ the radiance, $E_0$ the sun's irradiance at that wavelength and $\theta_s$ the solar zenith angle. The
$\pi$ converts radiance (per steradian) to the flux leaving a surface that scatters equally in all directions, and
$E_0 \cos\theta_s$ is the sunlight landing on a horizontal square metre when the sun is $\theta_s$ from overhead.

\begin{worked}[label={we:benchmarks:toa}]{Top-of-atmosphere reflectance}
Take round numbers\tagA{}: $L = 10$, $E_0 = 186$ (roughly the sun at 550 nm in the same units), $\theta_s = 30°$.
\begin{align*}
\pi \times 10 &= 31.42 && \text{numerator}\\
\cos 30° &= 0.866\\
186 \times 0.866 &= 161.1 && \text{denominator}\\
\rho_{\text{TOA}} &= 31.42 / 161.1 = 0.195.
\end{align*}
\end{worked}

That needs the solar zenith of every flight line, which the tiles do not carry, and it leaves the atmosphere in. The
preferred route - fetching AVIRIS-NG L2 reflectance for HAD100's ten flight lines - needs an Earthdata login and had
not been done when this was written.

\paragraph{The leakage check.} HAD100 comes from the same archive as the 79 AVIRIS-NG flight lines the models
trained on. If a benchmark scene were cut from a training line, a model could score well by memory. The README's
check\hfil\penalty50\hfilneg\tagM{research/experiments/benchmarks/README.md}:
\begin{itemize}
\item \textbf{No exact line overlap.} None of HAD100's 10 lines is among the 79.
\item \textbf{One near neighbour.} Line \texttt{ang20180227t082814}, near Kakinada in Andhra Pradesh, was flown on
the same day as three training lines; one of them, 18 minutes later, ran on the parallel track immediately east. The
outlines are about 630 m apart and do not overlap, and the four HAD100 tiles from that line sit 1,020--1,865 m from
the nearest training footprint.
\item Those 4 tiles become \textbf{6 of the 100 protocol scenes}: no pixel-level leakage, but near-domain proximity.
\textbf{Report HAD100 results with and without those 6 scenes.}
\end{itemize}
A second trap: HAD100's own training and test tiles come from the same 10 flight lines, so training on its anomaly-free
tiles would leak too.

\section{ABU}\label{sec:benchmarks:abu}

The Airport--Beach--Urban set \citep{kang2017hyperspectral} has 13 small scenes and is used by essentially every
hyperspectral anomaly paper. It has three problems for this product (README, section 4):
\begin{itemize}
\item twelve scenes are AVIRIS scaled at-sensor radiance with an unknown gain, not reflectance;
\item one scene, \texttt{beach-4}, is ROSIS-03 over Pavia, 430--860 nm only, min-max normalised to 0--1 over the whole
cube - not a physical quantity;
\item \textbf{no wavelengths}: the files hold \texttt{data} and \texttt{map} and nothing else, and which AVIRIS bands
were removed was never published.
\end{itemize}

Why ``radiance'', when nobody wrote it down? The median spectra peak in the blue at about 1,800 counts and fall to
about 20--100 near 2.4 $\mu$m\hfil\penalty50\hfilneg\tagM{research/experiments/benchmarks/README.md}:
\begin{align*}
20 / 1{,}800 &= 0.011, \text{ about } 1\%\\
100 / 1{,}800 &= 0.056, \text{ about } 6\%.
\end{align*}
Airports and cities do not reflect only 1--6\% as much at 2.4 $\mu$m as in the blue; sunlight itself is far weaker
there. A spectrum with the sun's shape still contains the sunlight: radiance. A model that reads bands by wavelength
cannot run on ABU without guessing them.

\paragraph{Protocols disagree too.} ``The first 50 bands'' means the first 50 of HAD100's 276 ``good'' bands
(452--697 nm) in one paper and the first 50 of the raw 425 (377--622 nm) in another. Both are visible-only; a model
using the short-wave infrared is handicapped against either. Quote the protocol and do not mix them in a table.

\paragraph{And satellites?} The implant-bench design (\file{docs/lld/implant-bench.md}) records the search: ``no
labelled PRISMA / EnMAP / Gaofen-5 anomaly set was found''. The four PRISMA scenes are the only labelled satellite
anomaly scenes the repository has.

\coursefig[0.9\linewidth]{ad-17-benchmarks}{What each benchmark supplies: the four things a model needs, plus
platform. Amber cells are the reason a benchmark cannot test this product as it is.}{fig:benchmarks:public}

So the options were: convert HAD100 to reflectance (a login, more work, still airborne), guess ABU's wavelengths (a
result resting on a guess), or put known targets into the repository's own held-out satellite scenes. The rest of
this chapter is the third.

\section{The implant idea, and what real objects look like}\label{sec:benchmarks:real}

An \index{implant bench}\emph{implant bench} takes real held-out scenes (\emph{hosts}) and adds synthetic objects
whose positions are known exactly. The pieces (\file{app/evaluation/implant.py} for the core,
\file{scripts/build_implant_bench.py} to build, \file{tests/test_evaluation/test_implant.py} for the hand-worked
tests). On 2026-10-17 the user asked for implanted targets ``somewhat similar to the ones we do have'', in ``subtle,
medium, easy'' categories. ``Similar'' needs measurement first, so \code{build_implant_bench.py real} measures every
labelled PRISMA object (4-connected; an object whose ring has fewer than 8 pixels is skipped)\hfil\penalty50\hfilneg\tagM{research/experiments/implant-bench/README.md}:

\begin{center}
\begin{tabular}{lc}
\toprule
 & 10th / 50th / 90th percentile \\
\midrule
size, pixels & 1 / 2 / 10 \\
single-pixel objects & 42\% \\
background NDVI & $-0.02$ / 0.37 / 0.60 \\
\bottomrule
\end{tabular}
\end{center}

Half of all real targets are one or two pixels. NDVI (normalised difference vegetation index) measures greenness;
the script's \code{ndvi_map} uses the mean of 450--650 nm as ``visible'' and of 800--900 nm as ``near-infrared'':
$\mathrm{NDVI} = (\mathrm{nir} - \mathrm{vis})/(\mathrm{nir} + \mathrm{vis})$.

\section{The ring and the contrast $z$}\label{sec:benchmarks:z}

``Subtle'' and ``easy'' need a number. The repository's is the contrast $z$, defined in the module docstring of
\file{app/evaluation/implant.py} (lines 5--11).

\begin{definition}[Ring]
The \index{ring}\emph{ring} of an object is the set of valid pixels 2 or 3 four-neighbour steps away from it, not
inside another object:
\begin{lstlisting}[basicstyle=\ttfamily\footnotesize,literate={~}{{\textasciitilde}}1]
# (H, W) bool: within 3 steps, minus within 1 step, valid only
ring = binary_dilation(obj, iterations=3) & ~binary_dilation(obj, iterations=1) & valid
\end{lstlisting}
\end{definition}

Each dilation step goes only up, down, left or right, so ``$k$ steps'' is \index{Manhattan distance}\emph{Manhattan
distance} $|\Delta y| + |\Delta x| = k$, and the ring is a diamond, not a square.

\paragraph{Counting the ring of a one-pixel object.} The pixels at Manhattan distance exactly $k \ge 1$ from a point
number $4k$: on each of the diamond's four sides there are $k$ pixels, counting from one corner up to (but not
including) the next. So
\begin{align*}
\text{distance 1: } & 4 \times 1 = 4 && \text{removed, too close}\\
\text{distance 2: } & 4 \times 2 = 8\\
\text{distance 3: } & 4 \times 3 = 12\\
\text{ring: } & 8 + 12 = 20.
\end{align*}
A square ring at the same distances would be a $7 \times 7$ block minus a $3 \times 3$ one, $49 - 9 = 40$ pixels.
The test \code{test_contrast_z_by_hand} pins 20, and the book's script agrees\tagC{}.

\coursefig[0.55\linewidth]{ad-18-ring}{The diamond ring around a one-pixel object. Distance 1 (grey) is skipped;
distances 2 (8 pixels) and 3 (12 pixels) form the 20-pixel ring (teal). The dashed square shows the 40-pixel ring a
$3 \times 3$ dilation would have made.}{fig:benchmarks:ring}

\begin{definition}[Contrast $z$]
Let $o_b$ be the object's mean reflectance in band $b$, and $r_b$, $\sigma_b$ the mean and the population standard
deviation of the ring in band $b$. Over the $B$ bands valid in both,
\begin{equation}\label{eq:benchmarks:z}
z = \frac{\frac1B\sum_b |o_b - r_b|}{\frac1B\sum_b \sigma_b}.
\end{equation}
\end{definition}

It is the standard-score idea of \cref{ch:rx}, averaged over bands: how far the object sits from its neighbourhood, in
units of the neighbourhood's own variation. The code adds $10^{-6}$ to the denominator so a perfectly flat ring does
not divide by zero.

\begin{worked}[label={we:benchmarks:z}]{The test's $z$, by hand}
Two bands, one-pixel object. Band 0 is 0.1 on the 8 distance-2 pixels and 0.3 on the 12 distance-3 pixels; band 1 is
0.5 everywhere. The object reads $(0.42, 0.5)$.
\begin{align*}
r_0 &= \frac{8 \times 0.1 + 12 \times 0.3}{20} = \frac{0.8 + 3.6}{20} = \frac{4.4}{20} = 0.22 \quad \text{(ring mean)}\\
\text{deviations} &: 0.1 - 0.22 = -0.12 \ (8\times), \quad 0.3 - 0.22 = +0.08 \ (12\times)\\
\sigma_0^2 &= \frac{8 \times 0.0144 + 12 \times 0.0064}{20} = \frac{0.1152 + 0.0768}{20} \quad \text{(divide by 20, not 19)}\\
&= \frac{0.192}{20} = 0.0096\\
\sigma_0 &= \sqrt{0.0096} = 0.0980\\
r_1 &= 0.5, \quad \sigma_1 = 0\\
|o_0 - r_0| &= |0.42 - 0.22| = 0.20, \quad |o_1 - r_1| = 0\\
z &= \frac{(0.20 + 0)/2}{(0.0980 + 0)/2} = \frac{0.10}{0.0490} = 2.04.
\end{align*}
The repository's \code{contrast_z} returns 2.0412\tagC{}.
\end{worked}

\paragraph{Tiers.} On the 909 real objects $z$ has median 0.68. Its terciles cut the tiers
(\file{app/evaluation/implant.py}, line 25):
\begin{lstlisting}[basicstyle=\ttfamily\footnotesize]
TIER_EDGES = (0.502, 0.934)      # z terciles of the 909 real PRISMA bench objects (2026-10-17)
def tier_of(z): return TIERS[int(z >= 0.502) + int(z >= 0.934)]   # subtle / medium / easy
\end{lstlisting}
By construction a third of real objects falls in each \index{tier}\emph{tier}. ``Subtle'' means ``as hard as the
hardest third of what analysts labelled'', not an absolute difficulty. The object of \cref{we:benchmarks:z}, at
$z = 2.04 \ge 0.934$, is easy.

\begin{keypoint}
Real and implanted objects must be measured with the same ring and the same bands, or their $z$ values are not
comparable and the tiers mean nothing.
\end{keypoint}

\section{Linear mixing: what an implant does to a pixel}\label{sec:benchmarks:mixing}

\paragraph{Derivation.} A sensor pixel records the average light from its ground footprint. Suppose a fraction $a$
of that footprint is covered by a material of reflectance spectrum $\vect{t}$, and the rest by the original background
of spectrum $\vect{b}$. The light from each part is proportional to its area times its reflectance, and the pixel
reports the area-weighted average:
\begin{align*}
\vect{x} &= \frac{a \cdot \vect{t} + (1 - a) \cdot \vect{b}}{a + (1 - a)} && \text{area-weighted average}\\
&= (1 - a)\,\vect{b} + a\,\vect{t}. && \text{the areas sum to 1}
\end{align*}
This is the \index{linear mixing}\emph{linear mixing model} (the sub-pixel mix of \cref{ch:material}), and $a$ is the
\index{abundance}\emph{abundance} or fill fraction. bench-v1's material implant is exactly this,
\code{implant_material}: \code{x' = (1 - f) * x + f * t}.

\paragraph{The mixing line.} Rearrange to see the geometry:
\begin{align*}
\vect{x} &= \vect{b} - a\,\vect{b} + a\,\vect{t} && \text{expand}\\
&= \vect{b} + a\,(\vect{t} - \vect{b}). && \text{collect the $a$ terms}
\end{align*}
As $a$ runs from 0 to 1, $\vect{x}$ moves along the straight segment from $\vect{b}$ to $\vect{t}$: the
\index{mixing line}\emph{mixing line}. The pixel's difference from the background is $\vect{x} - \vect{b} =
a(\vect{t} - \vect{b})$, linear in $a$, so a contrast measured against an unchanged background grows in proportion to
$a$. That is why one number, the strength, is enough to hit a target $z$.

\begin{worked}[label={we:benchmarks:mix}]{A pixel on the mixing line}
Two bands. Background $\vect{b} = (0.30, 0.20)$, material $\vect{t} = (0.10, 0.60)$\tagA{}. Then
$\vect{t} - \vect{b} = (-0.20, +0.40)$ and:
\begin{align*}
a = 0.25: \quad \vect{x} &= (0.30, 0.20) + 0.25 \times (-0.20, 0.40)\\
&= (0.30 - 0.05,\ 0.20 + 0.10) = (0.25, 0.30)\\
a = 0.50: \quad \vect{x} &= (0.30, 0.20) + 0.50 \times (-0.20, 0.40)\\
&= (0.30 - 0.10,\ 0.20 + 0.20) = (0.20, 0.40)
\end{align*}
Check with the first form at $a = 0.25$: band 0 is $0.75 \times 0.30 + 0.25 \times 0.10 = 0.225 + 0.025 = 0.25$. If
the ring around this pixel has $\sigma = (0.05, 0.05)$ and mean $\vect{b}$, then at $a = 0.25$:
$z = \frac{(0.05 + 0.10)/2}{0.05} = \frac{0.075}{0.05} = 1.5$, and at $a = 0.5$, $z = 3.0$: doubling $a$ doubles $z$.
\end{worked}

\begin{figure}[htbp]
\centering
\begin{tikzpicture}
\begin{axis}[width=7.5cm,height=7.5cm,xmin=0,xmax=0.4,ymin=0,ymax=0.7,xlabel={band 0 reflectance},
  ylabel={band 1 reflectance},grid=major,grid style={gray!20}]
\addplot[spec,very thick] coordinates {(0.30,0.20) (0.10,0.60)};
\addplot[only marks,mark=*,mark size=2.5pt,maskgrey] coordinates {(0.30,0.20)};
\addplot[only marks,mark=*,mark size=2.5pt,loss] coordinates {(0.10,0.60)};
\addplot[only marks,mark=o,mark size=2.5pt,spat,thick] coordinates {(0.25,0.30) (0.20,0.40)};
\node[anchor=west,font=\small] at (axis cs:0.305,0.20) {$\vect{b}$ (a = 0)};
\node[anchor=west,font=\small,text=loss] at (axis cs:0.105,0.62) {$\vect{t}$ (a = 1)};
\node[anchor=west,font=\small,text=spat] at (axis cs:0.255,0.31) {$a = 0.25$};
\node[anchor=west,font=\small,text=spat] at (axis cs:0.205,0.41) {$a = 0.5$};
\end{axis}
\end{tikzpicture}
\caption{The mixing line of \cref{we:benchmarks:mix} in the plane of two bands: every mixture of background
$\vect{b}$ and material $\vect{t}$ lies on the straight segment between them, at a distance from $\vect{b}$
proportional to the abundance $a$.}
\label{fig:benchmarks:mixline}
\end{figure}

\paragraph{The signature implant.} bench-v1's second source is additive: \code{implant_signature} sets
$\vect{x}' = \max(\vect{x} + a\,\vect{d}, 0)$, where $\vect{d} = \vect{o} - \vect{r}$ is a real object's difference
from its ring, carried to the host's wavelengths by interpolation. The $\max(\cdot, 0)$ stops reflectance going
negative. Both implants move the object's mean spectrum along a straight line: if $\vect{o}_0$ is the object's mean
before implanting,
\begin{equation}\label{eq:benchmarks:line}
\vect{o}(s) = \vect{o}_0 + s\,\vect{\Delta}, \qquad
\vect{\Delta} = \begin{cases} \vect{d} & \text{signature, } s = a\\ \vect{t} - \vect{o}_0 & \text{material, } s = f \le 1. \end{cases}
\end{equation}
For the material, $(1 - f)\vect{o}_0 + f\vect{t} = \vect{o}_0 + f(\vect{t} - \vect{o}_0)$, the mixing line again.

\begin{sourcenote}[: signatures and descending wavelengths]
Until commit b94dbfd every signature implant was garbled: PRISMA lists its bands in descending wavelength, and
\code{np.interp} needs ascending sample points. \code{interp_spectrum} now sorts first (\file{app/evaluation/implant.py},
lines 160--166). \Cref{ch:sortbug} tells the whole story; every signature-implant number before 2026-10-18 is void.
\end{sourcenote}

\section{Solving for the strength on the realised contrast}\label{sec:benchmarks:solve}

Each object gets a target $z$, drawn from the real objects of its tier. Put \cref{eq:benchmarks:line} into
\cref{eq:benchmarks:z}, with the ring unchanged:
\begin{equation}\label{eq:benchmarks:zs}
z(s) = \frac{\frac1B\sum_b |o_{0,b} + s\,\Delta_b - r_b|}{\bar\sigma}, \qquad \bar\sigma = \frac1B\sum_b \sigma_b.
\end{equation}
A first version of the design gave the simpler $s = z\,\bar\sigma / \overline{|\Delta|}$. That assumes $\vect{o}_0 =
\vect{r}$: that the spot looked exactly like its ring before anything was added. Real pixels never do.

\begin{sourcenote}[: \texttt{strength\_for\_z} never existed]
The implant-bench LLD's first interface list named \code{strength_for_z} with the simple formula. That function was
never in the source; the code has \code{solve_strength}, which solves \cref{eq:benchmarks:zs}
(\file{app/evaluation/implant.py}, lines 80--104). The LLD now lists \code{solve_strength}.
\end{sourcenote}

\begin{worked}[label={we:benchmarks:solve}]{One band: the simple formula misses, the realised solve hits}
Ring mean $r = 0.20$, $\sigma = 0.05$; the spot reads $o_0 = 0.22$, already 0.02 above its ring; the implant darkens,
$\Delta = -0.10$; target $z = 1.0$.
\begin{align*}
\text{simple: } s &= 1.0 \times 0.05 / 0.10 = 0.5\\
\text{realised at } s = 0.5: \ z &= |0.22 + 0.5 \times (-0.10) - 0.20| / 0.05\\
&= |{-0.03}| / 0.05 = 0.6 \quad \text{(medium, not easy)}
\end{align*}
The natural $+0.02$ cancelled part of the darkening. Solve \cref{eq:benchmarks:zs} instead:
\begin{align*}
z(s) &= |0.02 - 0.10\,s| / 0.05 && o_0 - r = 0.02\\
z(0) &= 0.02 / 0.05 = 0.4 && \text{the spot's own contrast}\\
z(s) &= 0 \text{ at } s = 0.2 && \text{implant cancels the offset}\\
0.02 - 0.10\,s &= -0.05 && \text{rising branch, } |\cdot| = 0.05\\
s &= 0.07 / 0.10 = 0.7\\
\text{check: } z(0.7) &= |0.22 - 0.07 - 0.20| / 0.05 = 0.05 / 0.05 = 1.0.
\end{align*}
The repository's \code{solve_strength} returns $s = 0.7$, $z = 1.0$; with $\Delta = +0.10$ it returns $s = 0.3$\tagC{}.
\end{worked}

\paragraph{Why the first crossing after the minimum.} Each term $|o_{0,b} + s\Delta_b - r_b|$ is the absolute value of
a straight line in $s$: a ``V'' shape, which is \emph{convex} (the segment between any two points of its graph lies on
or above the graph). A sum of convex functions is convex, and dividing by the positive $\bar\sigma$ keeps it convex.
A convex function of one variable falls (or stays flat) up to its minimum and rises from then on. So $z(s)$ can dip
below $z(0)$ and then climb for good; the useful answer is the first $s$ at or after the minimum where $z(s)$ reaches
the target. The code evaluates $z$ on a grid of 400 strengths, finds the minimum, takes the first grid point at or
above the target to its right, and refines by linear interpolation between that point and the one before.

Two more rules follow from the shape:
\begin{enumerate}
\item If $z(0) > z_{\text{target}}$, the spot is already too contrasty; \code{solve_strength} returns \code{None} and
the builder tries another spot (up to 200 times).
\item A material may not reach the target even at full fill $f = 1$; then $f$ is capped and the lower realised $z$ is
recorded (2\% of bench-v1's material implants\hfil\penalty50\hfilneg\tagM{research/experiments/implant-bench/README.md}).
\end{enumerate}
After implanting, $z$ is measured again, and \textbf{the realised $z$ and its tier are what the benchmark records}.
One consequence: a subtle target of $z = 0.2$ can only go on a spot whose own $z(0) \le 0.2$, so subtle implants sit
on the most uniform spots of a scene.

\coursefig[0.7\linewidth]{ad-18-strength}{$z(s)$ for \cref{we:benchmarks:solve}. The simple formula's $s = 0.5$ lands at
$z = 0.6$ (medium); the solver follows the curve through its minimum at $s = 0.2$ and crosses the target at
$s = 0.7$.}{fig:benchmarks:strength}

\section{bench-v1, and how it failed}\label{sec:benchmarks:v1}

bench-v1 was built on 2026-10-17 with seed 0: six held-out test scenes per sensor as hosts, objects placed on valid,
cloud- and water-free land at the real object density, with centres weighted so their background NDVI mix matches the
real objects'\hfil\penalty50\hfilneg\tagM{research/experiments/implant-bench/README.md}.

\begin{center}
\begin{tabular}{lccc}
\toprule
sensor & hosts & objects & subtle / medium / easy \\
\midrule
EnMAP (224 bands) & 6 & 1,179 & 385 / 412 / 382 \\
PRISMA (234 bands) & 6 & 1,318 & 429 / 463 / 426 \\
AVIRIS-NG (372 bands) & 6 & 839 & 273 / 295 / 271 \\
total & 18 & 3,336 & 1,087 / 1,170 / 1,079 \\
\bottomrule
\end{tabular}
\end{center}

Size and contrast matched the real objects closely (size percentiles 1 / 2 / 9 against 1 / 2 / 10; per-tier $z$
percentiles within a few hundredths); greenness did not (ring NDVI median 0.49 against 0.37, no host had the bright
bare ground of real scene set4). 97\% of objects landed in their targeted tier.

Then the stage-3 evaluation (\cref{ch:runeval}) ran a model-free detector beside the models. PCA-RX, which is local RX
on the cube's own top 30 principal components, caught 30.2\% of bench-v1 objects at the top 1\% against the best
model's 26.9\%, and lost by a wide margin on the real objects (pAUC 0.128 against 0.158)\hfil\penalty50\hfilneg\tagM{docs/lld/evaluation.md}. A
benchmark on which a 1990s detector beats every model, while real labels say the opposite, is measuring something
else. A diagnosis on the PRISMA hosts found two differences between implants and real objects: foreignness and halo.

\section{Foreignness}\label{sec:benchmarks:foreign}

Start from an object's difference from its ring, $\vect{d} = \vect{o} - \vect{r}$ (a vector with one entry per band).

\begin{definition}[Background subspace and foreign fraction]
The \index{background subspace}\emph{background subspace} of a scene is spanned by its top $K = 10$ principal
directions (\code{FOREIGN_K}), the orthonormal rows of a $K \times B$ matrix $\mat{P}$ computed by
\code{background_subspace} (\cref{ch:linalg} for principal directions). The part of $\vect{d}$ inside the subspace is
its projection $\vect{d}_{\text{in}} = \mat{P}\T\mat{P}\vect{d}$; the rest, $\vect{d}_{\text{out}} = \vect{d} -
\vect{d}_{\text{in}}$, is outside. The \index{foreign fraction}\emph{foreign fraction} is the share of $\norm{\vect{d}}^2$
outside:
\begin{equation}\label{eq:benchmarks:foreign}
\phi(\vect{d}) = \frac{\norm{\vect{d}_{\text{out}}}^2}{\norm{\vect{d}}^2} = 1 - \frac{\norm{\mat{P}\vect{d}}^2}{\norm{\vect{d}}^2}.
\end{equation}
\end{definition}

The second form is what \code{foreign_fraction} computes. It holds because the inside and outside parts are
perpendicular, so by Pythagoras $\norm{\vect{d}}^2 = \norm{\vect{d}_{\text{in}}}^2 + \norm{\vect{d}_{\text{out}}}^2$,
and because $\mat{P}$ has orthonormal rows, $\norm{\vect{d}_{\text{in}}} = \norm{\mat{P}\T\mat{P}\vect{d}} =
\norm{\mat{P}\vect{d}}$.

\begin{worked}[label={we:benchmarks:foreign}]{Foreignness of $(3, 4, 0)$}
Three bands; the scene's single main direction is band 0, so $\mat{P} = (1, 0, 0)$. Take $\vect{d} = (3, 4, 0)$.
\begin{align*}
\norm{\vect{d}}^2 &= 3^2 + 4^2 + 0^2 = 9 + 16 + 0 = 25\\
\mat{P}\vect{d} &= 1 \times 3 + 0 \times 4 + 0 \times 0 = 3, \quad \norm{\mat{P}\vect{d}}^2 = 9 && \text{inside: } (3, 0, 0)\\
\vect{d}_{\text{out}} &= (3, 4, 0) - (3, 0, 0) = (0, 4, 0), \quad \norm{\vect{d}_{\text{out}}}^2 = 16\\
\phi &= 16 / 25 = 0.64 = 1 - 9/25.
\end{align*}
\end{worked}

The diagnosis medians on PRISMA\hfil\penalty50\hfilneg\tagM{docs/training/anomaly-detection/19-running-an-evaluation.md}: real objects
\textbf{0.006}, bench-v1 implants \textbf{0.042}, seven times as foreign. Real objects are unusual \emph{mixtures of
the scene's own materials}; bench-v1's implants pointed in directions the scene never takes. Whitening (\cref{ch:rx})
rewards exactly that: a difference of 0.01 along a direction where the scene varies with spread 0.05 is $0.2$ spreads
(squared, 0.04); along a direction of spread 0.002 it is 5 spreads (squared, 25), $25/0.04 = 625$ times the score.

\paragraph{Hitting a target foreignness.} \code{match_foreignness} keeps $\vect{d}_{\text{in}}$ and rescales
$\vect{d}_{\text{out}}$ by a factor $\beta$, so that $\vect{d}' = \vect{d}_{\text{in}} + \beta\,\vect{d}_{\text{out}}$
has foreign fraction $t$. Derive $\beta$. Write $A = \norm{\vect{d}_{\text{in}}}^2$, $C = \norm{\vect{d}_{\text{out}}}^2$:
\begin{align*}
\phi(\vect{d}') &= \frac{\beta^2 C}{A + \beta^2 C} = t && \text{parts stay perpendicular}\\
\beta^2 C &= t A + t \beta^2 C && \text{multiply both sides by } A + \beta^2 C\\
\beta^2 C (1 - t) &= t A && \text{collect the } \beta^2 \text{ terms}\\
\beta^2 &= \frac{t}{1 - t} \cdot \frac{A}{C}. && \text{divide by } C(1 - t)
\end{align*}
This is the formula in the docstring of \code{match_foreignness} (lines 137--148).

\begin{worked}[label={we:benchmarks:match}]{Matching $(4, 0, 2)$ to foreignness 0.1}
Same $\mat{P} = (1, 0, 0)$; $\vect{d} = (4, 0, 2)$, so $\vect{d}_{\text{in}} = (4, 0, 0)$, $A = 16$, and
$\vect{d}_{\text{out}} = (0, 0, 2)$, $C = 4$.
\begin{align*}
t = 0.2: \quad \beta^2 &= \tfrac{0.2}{0.8} \times \tfrac{16}{4} = 0.25 \times 4 = 1, \quad \beta = 1 && \text{already at target}\\
t = 0.1: \quad \beta^2 &= \tfrac{0.1}{0.9} \times 4 = 0.4444, \quad \beta = 0.6667\\
\vect{d}' &= (4, 0, 0) + 0.6667 \times (0, 0, 2) = (4, 0, 1.3333)\\
\text{check: } \phi &= \frac{1.3333^2}{16 + 1.3333^2} = \frac{1.7778}{17.7778} = 0.1.
\end{align*}
The repository's function returns $(4, 0, 1.3333)$\tagC{}.
\end{worked}

\section{The halo and the Gaussian point-spread function}\label{sec:benchmarks:psf}

\begin{definition}[Halo and halo ratio]
The \index{halo}\emph{halo} of an object is the ring of valid pixels touching it (Manhattan distance 1). Its
\emph{halo ratio} is the halo's contrast $z$ (against the same distance-2--3 ring) over the object's.
\end{definition}

A ratio of 0 is a hard edge; near 1 the object bleeds fully into its neighbours. Real sensors blur - each pixel collects
some light from next door - so real objects have halos: the diagnosis measured a median halo ratio of 0.74 for real
objects against 0.60 for bench-v1 implants (diagnosis ruler; \cref{sec:benchmarks:v2} explains the re-measured
0.643).

\begin{definition}[Gaussian point-spread function]
A \index{point-spread function}\emph{point-spread function} (PSF) is the shape a single bright point takes in the
image. A \emph{Gaussian} PSF of standard deviation $\sigma$ pixels spreads a point at the origin so that the pixel at
offset $(\Delta y, \Delta x)$ receives a share proportional to
\begin{equation}\label{eq:benchmarks:gauss}
g(\Delta y, \Delta x) = \exp\!\Big(-\frac{\Delta y^2 + \Delta x^2}{2\sigma^2}\Big).
\end{equation}
\end{definition}

bench-v2 adds its signal through a blurred footprint. \code{psf_weights} (lines 151--157) blurs the object mask with
\code{ndimage.gaussian_filter} and divides by the blurred mask's mean over the object, so the object keeps weight 1 on
average and its neighbours get the tail:
\begin{lstlisting}[basicstyle=\ttfamily\footnotesize]
w = ndimage.gaussian_filter(obj.astype(np.float64), sigma, mode="constant")   # (H, W)
return w / max(w[obj].mean(), 1e-12)                                          # mean 1 on obj
\end{lstlisting}
The build uses $\sigma = 1.4$ (\texttt{-{}-psf-sigma 1.4}, in \file{research/experiments/implant-bench/README.md} and
the LLD's build lines), chosen by calibration builds rather than from a formula. For a one-pixel object the blurred
mask is the Gaussian kernel itself, so after dividing by the centre value the weight at $(\Delta y, \Delta x)$ is
exactly \cref{eq:benchmarks:gauss}. With $\sigma = 1.4$, $2\sigma^2 = 2 \times 1.96 = 3.92$.

\begin{worked}[label={we:benchmarks:psf}]{The PSF weights and the abundance map on a $3 \times 3$ footprint}
Do this one on paper with a calculator. One-pixel object at the centre of a $3 \times 3$ window.
\begin{align*}
\text{centre } (0,0): & \quad \exp(0) = 1\\
\text{edge } (\pm1, 0), (0, \pm1): & \quad \exp(-1/3.92) = \exp(-0.2551) = 0.7748\\
\text{corner } (\pm1, \pm1): & \quad \exp(-2/3.92) = \exp(-0.5102) = 0.6004
\end{align*}
Now the abundance. A bench-v2 material implant uses $\vect{d}_0 = \vect{t} - \vect{r}$ (library spectrum minus ring
mean, \file{scripts/build_implant_bench.py}, line 238) and adds $s\,w(p)\,\vect{d}$ to each pixel $p$. Where the pixel
looks like its ring, $\vect{x} \approx \vect{r}$:
\begin{align*}
\vect{x}' &= \vect{x} + s\,w(p)\,(\vect{t} - \vect{r}) && \text{the v2 implant}\\
&\approx \vect{x} + s\,w(p)\,(\vect{t} - \vect{x}) && \text{use } \vect{x} \approx \vect{r}\\
&= (1 - a(p))\,\vect{x} + a(p)\,\vect{t}, \quad a(p) = s\,w(p). && \text{the mixing model}
\end{align*}
So the PSF turns one strength into an \emph{abundance map}. At strength $s = 0.1$\tagA{}:
\[
a = 0.1 \times \begin{pmatrix} 0.6004 & 0.7748 & 0.6004\\ 0.7748 & 1 & 0.7748\\ 0.6004 & 0.7748 & 0.6004\end{pmatrix}
= \begin{pmatrix} 0.0600 & 0.0775 & 0.0600\\ 0.0775 & 0.1000 & 0.0775\\ 0.0600 & 0.0775 & 0.0600\end{pmatrix}.
\]
The centre pixel is 10\% material; each edge neighbour 7.75\%; each corner 6\%. The script's \code{psf_weights} gives the
same matrix\tagC{}.
\end{worked}

\begin{figure}[htbp]
\centering
\begin{tikzpicture}[x=1.15cm,y=1.15cm]
\foreach \y/\r in {0/{0.0600,0.0775,0.0600},1/{0.0775,0.1000,0.0775},2/{0.0600,0.0775,0.0600}}{
  \foreach \v [count=\x from 0] in \r {
    \pgfmathsetmacro{\shade}{\v*900}
    \fill[loss!\shade!white] (\x,\y) rectangle ++(1,1);
    \draw[gray!60] (\x,\y) rectangle ++(1,1);
    \node[font=\small] at (\x+0.5,\y+0.5) {\v};
  }
}
\draw[spat,very thick] (1,1) rectangle ++(1,1);
\node[anchor=west,font=\small,align=left] at (3.3,1.5) {object pixel (indigo box): $a = s = 0.1$\\
  edge neighbours: $a = 0.1 \times 0.7748$\\ corners: $a = 0.1 \times 0.6004$\\
  halo = the four edge cells};
\end{tikzpicture}
\caption{The abundance map of \cref{we:benchmarks:psf}: a one-pixel bench-v2 implant at strength 0.1 with a Gaussian
PSF of $\sigma = 1.4$, shown on its $3 \times 3$ footprint (darker amber = more material). The blur continues
beyond the window, weaker with distance.}
\label{fig:benchmarks:abund}
\end{figure}

\paragraph{The ring is lifted too: the shrink factor.} The PSF tail reaches the ring. For a one-pixel object, average
\cref{eq:benchmarks:gauss} over the 20 ring cells (distance 2: 4 cells at $(\pm2,0),(0,\pm2)$ and 4 at
$(\pm1,\pm1)$; distance 3: 4 at $(\pm3,0),(0,\pm3)$ and 8 at $(\pm2,\pm1),(\pm1,\pm2)$):
\begin{align*}
4 \times e^{-4/3.92} &= 4 \times 0.3604 = 1.4418\\
4 \times e^{-2/3.92} &= 4 \times 0.6004 = 2.4015\\
4 \times e^{-9/3.92} &= 4 \times 0.1007 = 0.4027\\
8 \times e^{-5/3.92} &= 8 \times 0.2793 = 2.2343\\
\bar w_{\text{ring}} &= (1.4418 + 2.4015 + 0.4027 + 2.2343) / 20 = 6.4803 / 20 = 0.3240.
\end{align*}
The script's \code{psf_weights} gives 0.3240\tagC{}. The object mean moves by $s \times 1 \times \vect{d}$; the ring mean
by $s \times 0.3240 \times \vect{d}$; so the object's difference from its ring moves by only
\begin{equation}\label{eq:benchmarks:shrink}
s\,(1 - \bar w_{\text{ring}})\,\vect{d} = 0.676\,s\,\vect{d}.
\end{equation}
The builder passes this \code{shrink} $= 1 - \bar w_{\text{ring}}$ into the strength solve.

\paragraph{The halo ratio a PSF produces, on flat ground.} With a perfectly flat background the halo and the object
are measured against the same ring, so the ring's $\bar\sigma$ cancels in the ratio of their $z$ values:
\begin{align*}
\text{halo ratio} &= \frac{|w_{\text{halo}} - \bar w_{\text{ring}}|\; s\,\overline{|\vect{d}|}/\bar\sigma}
                         {|1 - \bar w_{\text{ring}}|\; s\,\overline{|\vect{d}|}/\bar\sigma}
  && \text{both differences scale } \vect{d}\\
&= \frac{0.7748 - 0.3240}{1 - 0.3240} = \frac{0.4508}{0.6760} = 0.667. && s,\ \vect{d},\ \bar\sigma \text{ cancel}
\end{align*}
Close to the real 0.643 measured on the bench's own ruler. Real halos also carry the background's texture, which is
why $\sigma$ was calibrated on builds, not fixed by this formula.

\section{bench-v2: solve on the object as it ends up}\label{sec:benchmarks:v2}

bench-v2 (2026-10-17) copies three measured properties from real objects: $z$, foreign fraction and halo ratio, with
$z$ and foreignness drawn \emph{as a pair from one real object} of the chosen tier. They are related: subtle real
objects are about ten times as foreign as easy ones (median 0.0238 against 0.0024 on \file{real_v2.json}'s
ruler\hfil\penalty50\hfilneg\tagM{docs/training/anomaly-detection/20-bench-v2.md}), so drawing them independently would create objects
that do not exist. A third source joined signature and material: a \emph{host} spectrum taken from a land pixel at
least 100 pixels away in the same scene. The first calibration build missed, and three bugs were found.

\begin{enumerate}
\item \textbf{Measured on the wrong bands.} Foreignness and $z$ were computed over every host band, including the noisy
water-absorption and detector-edge bands the evaluator never reads. Now \code{evaluation_bands} keeps the product's
bands only (exclusions removed, and only bands at least 60\% finite). Measure a benchmark with the ruler that will
grade on it.
\item \textbf{One gap deleted a band.} \code{background_subspace} centres pixels by subtracting each band's mean. A
plain mean with one missing value is missing: $(0.20 + \text{NaN} + 0.40)/3 = \text{NaN}$, and the whole band dropped
out. \code{np.nanmean} gives $(0.20 + 0.40)/2 = 0.30$; centred values $-0.10$, gap, $+0.10$; the gap is then set to 0,
``the band mean''. 130--182 of 234 PRISMA host bands had at least one gap (docstring, lines 115--118).
\item \textbf{Solved on the signal, not on the object.} The first version matched the foreignness of the added signal.
But the evaluator sees the realised difference
\begin{equation}\label{eq:benchmarks:realised}
\vect{o} - \vect{r} = (\vect{o}_0 - \vect{r}) + s \cdot \text{shrink} \cdot \vect{d}',
\end{equation}
and the spot's natural offset $\vect{o}_0 - \vect{r}$ is already foreign (0.009--0.045 on random un-implanted spots).
\end{enumerate}

\begin{worked}[label={we:benchmarks:realised}]{An unforeign signal on a foreign spot}
Two bands; the scene's main direction is band 0. Natural offset $\vect{o}_0 - \vect{r} = (0.00, 0.01)$, wholly outside.
Signal $s \cdot \text{shrink} \cdot \vect{d}' = (0.05, 0.00)$, wholly inside: its foreignness is 0.
\begin{align*}
\vect{o} - \vect{r} &= (0.00 + 0.05,\ 0.01 + 0.00) = (0.05, 0.01)\\
\norm{\cdot}^2 &= 0.05^2 + 0.01^2 = 0.0025 + 0.0001 = 0.0026, \quad \text{outside } 0.0001\\
\phi &= 0.0001 / 0.0026 = 0.038.
\end{align*}
The signal was perfectly unforeign; the object the evaluator sees is five times the real median (0.0071).
\end{worked}

\code{solve_realised} (lines 172--187) fixes it. It tries a ladder of signal foreignness levels,
\code{FOREIGN_LADDER} $= (0, 0.0005, 0.001, \ldots, 0.25, 0.5)$; at each it solves the strength for the target $z$
with \code{solve_strength} on \code{shrink * d}, measures the foreignness of the realised difference
\cref{eq:benchmarks:realised}, and keeps the level whose realised foreignness is closest to the target on a log scale,
$|\ln(\phi_{\text{realised}} / \phi_{\text{target}})|$. A spot whose own offset is so foreign that no level comes
within a factor \code{tol} of the target (1.5 in bench-v2, 1.2 in bench-v2.1) is rejected and another spot tried.

\begin{keypoint}
Match a benchmark property on the object the evaluator will see, never on the ingredient you added.
\end{keypoint}

\section{Accepting a benchmark: the ordering test}\label{sec:benchmarks:accept}

bench-v2 had two acceptance tests, written before the build.

\paragraph{Acceptance 1 - does it look like the real objects?} Foreignness and halo medians within 25\% of the real
objects', overall and in every tier. The calibration build on 2 PRISMA hosts passed: foreignness 0.0071 real against
0.0081, halo 0.643 against 0.642, $z$ 0.69 against 0.74, every tier within 0.90--1.15$\times$\hfil\penalty50\hfilneg\tagM{research/experiments/implant-bench/README.md}.
The final bench, measured on all 6 PRISMA hosts, came out at 1.10 / 1.28 / 1.27$\times$ real foreignness (subtle /
medium / easy) - marginally over the limit in two tiers - and halo 0.93--1.07$\times$.

\paragraph{Acceptance 2 - does it reproduce the real ordering?} The real bench says the fine-tuned FT03 beats PCA-RX
(pAUC 0.158 against 0.128). bench-v2 may judge models only if it says the same.

\begin{drift}
bench-v2, \file{configs/eval/bench_v2.json}, 18 hosts, 2,913 objects, residual Mahalanobis, 2026-10-17
\hfil\penalty50\hfilneg\tagM{research/experiments/implant-bench/README.md}. Chance: 10.1\% / 1.1\% / 0.005.

\begin{center}
\small
\begin{tabular}{lccccc}
\toprule
 & caught top 1\% (0.1\%) & pAUC & easy & medium & subtle \\
\midrule
F001 ep109 & 29.2\% (10.2\%) & 0.145 & 48.5\% & 23.1\% & 13.1\% \\
FT02c ep10 & 32.8\% (11.4\%) & 0.166 & 52.9\% & 27.6\% & 14.7\% \\
FT03 ep10 & 32.3\% (11.4\%) & 0.164 & 52.7\% & 26.9\% & 14.2\% \\
PCA-RX (no model) & 30.2\% (9.1\%) & 0.154 & 47.7\% & 24.9\% & 15.2\% \\
\bottomrule
\end{tabular}
\end{center}
(v0.2.0's row, 22.1\%, is void: a reading artefact of the sort bug, \cref{ch:sortbug}.) In objects at the top 1\%:
FT03 $0.323 \times 2{,}913 = 940.9$, about 941; PCA-RX $0.302 \times 2{,}913 = 879.7$, about 880.
\end{drift}

Met: both fine-tuned models above PCA-RX on every headline measure (subtle is a tie). And the subtle tier, at chance on
bench-v1, is now above it (13--15\% against 10.1\%).

\coursefig[0.85\linewidth]{ad-20-three-benches}{pAUC to 1\% false alarms for FT03 (amber) and PCA-RX (teal) on three
benchmarks. The real PRISMA bench puts FT03 first; bench-v1 reversed that order; bench-v2 restores
it.}{fig:benchmarks:three}

\coursefig[0.85\linewidth]{iw-57-what-v2-must-match}{What bench-v2 had to match: medians on PRISMA, real objects
(grey) against bench-v1 implants (amber), overall and in the easy tier, on the diagnosis ruler; the shaded band is
within 25\% of the real value.}{fig:benchmarks:mustmatch}

\paragraph{Why the ordering test is the one that matters.} Acceptance 1 checks that the implants \emph{look} right on
the properties somebody thought to measure. Acceptance 2 checks that the benchmark \emph{behaves} right on a fact
nobody built into it. bench-v1 would have passed a looser acceptance 1 - its sizes and contrasts matched - and still
ranked a classical detector first. A benchmark is a model of reality; test it on something you already know before you
use it to learn something you do not.

\section{bench-v2.1}\label{sec:benchmarks:v21}

After the descending-wavelength fix (commit b94dbfd) every signature implant had to be rebuilt, and the foreignness
tolerance was tightened from 1.5 to 1.2: bench-v2.1, built 2026-10-18 with \texttt{-{}-version 2 -{}-psf-sigma 1.4
-{}-foreign-tol 1.2 -{}-seed 0}, 18 hosts, 2,817 objects (signature 985, material 927, host 905).

\begin{drift}
bench-v2.1, \file{configs/eval/bench_v2_1.json}, residual Mahalanobis, 2026-10-18
\hfil\penalty50\hfilneg\tagM{research/experiments/implant-bench/README.md}.
\begin{center}
\small
\begin{tabular}{lccccc}
\toprule
model & caught top 1\% (0.1\%) & pAUC & easy & medium & subtle \\
\midrule
v0.2.0 (read correctly) & 35.5\% (12.6\%) & 0.168 & 53.7\% & 29.1\% & 19.5\% \\
F001 ep109 & 33.1\% (11.7\%) & 0.157 & 51.7\% & 26.7\% & 16.9\% \\
FT02c ep10 & 35.1\% (12.8\%) & 0.169 & 53.7\% & 28.7\% & 18.6\% \\
FT03 ep10 & 34.9\% (12.7\%) & 0.168 & 53.7\% & 28.3\% & 18.6\% \\
FT05 ep5 & 35.4\% (12.8\%) & 0.170 & 54.3\% & 28.9\% & 18.8\% \\
PCA-RX & 33.2\% (9.8\%) & 0.154 & 50.4\% & 26.9\% & 18.4\% \\
\bottomrule
\end{tabular}
\end{center}
Per-tier foreignness against real (PRISMA hosts): 1.05 / 1.15 / \textbf{1.60}$\times$; the tighter tolerance skips
low-foreignness easy targets, so the survivors skew high. Halo 0.92--1.06$\times$.
\end{drift}

Read correctly, v0.2.0 ties or beats every new model here, as it ties them on the real bench; all models stay above
PCA-RX at the top 0.1\% (12.6--12.8\% against 9.8\%), as on the real bench; at the top 1\% pretraining alone (F001,
33.1\%) ties PCA-RX (33.2\%). Since 2026-10-17 vessels on water are in scope (\cref{ch:runeval}), but
\code{place_and_implant} still removes water and cloud from the eligible area, so no bench-v2.x object sits on water.

\begin{confusions}
\confusion{HAD100 is AVIRIS-NG and we train on AVIRIS-NG, so it fits?}{Same sensor, different product: the training
data is L2 reflectance, HAD100 ships L1B radiance.}
\confusion{No exact flight-line overlap means no leakage worry?}{Six scenes are 1--2 km from a training swath flown the
same day. Not pixel leakage, but not unseen ground either; report both ways.}
\confusion{The ring is the $3 \times 3$ neighbourhood?}{It is a 20-pixel diamond at Manhattan distance 2--3 for a
one-pixel object; distance 1 (the halo) is skipped.}
\confusion{The tier is the one the builder aimed for?}{It is the realised one, measured after implanting.}
\confusion{The foreignness of the added signal is the object's foreignness?}{Only on a spot identical to its ring. The
evaluator sees natural offset plus signal (\cref{we:benchmarks:realised}).}
\confusion{The PSF $\sigma$ and the halo ratio are the same number?}{The PSF sets the blur; the halo ratio is measured
afterwards. $\sigma = 1.4$ gives 0.667 on flat ground and about 0.64 on real texture.}
\confusion{bench-v2 passed, so its absolute numbers are what the product achieves?}{It passed a ranking test. Real
objects are caught at about 45\% at the top 1\%, bench-v2's at about 32\%. Compare models within one bench.}
\end{confusions}

\begin{checkyourself}
\item\label{ex:benchmarks:1} Count the ring of a two-pixel horizontal object (pixels at Manhattan distance 2 or 3 from
the nearer object pixel). Count row by row.
\item\label{ex:benchmarks:2} A ring has mean 0.30 and $\sigma$ 0.06 in each of two bands. An object reads 0.36 and
0.27. Compute $z$ and its tier.
\item\label{ex:benchmarks:3} One band: $r = 0.40$, $\sigma = 0.10$, $o_0 = 0.43$, $\Delta = -0.20$. What strength
gives $z = 1.5$? What does the simple formula give, and what $z$ does it realise?
\item\label{ex:benchmarks:4} With $\mat{P} = (1, 0, 0)$, match $\vect{d} = (3, 0, 3)$ to foreignness 0.1: find
$\beta$ and $\vect{d}'$, and verify.
\item\label{ex:benchmarks:5} Compute the PSF weight at offset $(2, 0)$ and $(1, 1)$ for $\sigma = 1.0$, the ring mean
weight of a one-pixel object at $\sigma = 1.0$ (use $e^{-2} = 0.1353$, $e^{-1} = 0.3679$, $e^{-4.5} = 0.0111$,
$e^{-2.5} = 0.0821$, $e^{-0.5} = 0.6065$), and the flat-ground halo ratio. Is the blur sharper or softer than the build's?
\item\label{ex:benchmarks:6} A spot's natural offset is $(0.00, 0.02)$; the added signal is $(0.06, 0.00)$; band 0 is
the scene's main direction. What is the realised foreignness, and what did the signal alone have?
\end{checkyourself}

\nextchapter{ch:runeval}
